\section{Parameters Description}

\subsection{Initial Parameters}

As mentioned in the introduction, our objective is to find such parameters that will be better than SLH-DSA. 
For the Bitcoin we require no more than 128-bit security. 
If we put the security cap higher this will lead to bigger signature size and costs, which is not preferable.
Another parameter that sets the further foundation for the parameter search is the number of signatures, namely $q_{s}$. 
In SLH-DSA-128s version, the number of signatures is set to be $q_{s} = 2^{64}$. 
While parameters for standardized SLH-DSA are created for general purpose usage, the number of required signatures in Bitcoin is lower.
Setting $q_{s} = 2^{40}$ is enough for our purposes, which will decrease the size and the computational cost.

\subsection{SLH-DSA Parameters}

The basic architecture of SLH-DSA relies on a hyper-tree of extended Merkle signature schemes (XMSS) to authenticate underlying one-time and few-time signature keys. 
We formally parameterize each candidate configuration using the tuple $(h, d, k, a, w)$, where each variable controls a distinct geometric trade-off within the overall structure.

\paragraph{Hyper-Tree Geometry ($h$ and $d$)}
The hyper-tree configuration determines the capacity and layer traversal costs of the signature scheme.
\begin{itemize}
    \item $h$ represents the total height of the multi-layered hyper-tree. 
    This variable dictates the maximum number of few-time signature instances supported at the base layer, yielding a total capacity of $2^h$ unique leaf nodes.
    \item $d$ defines the number of independent tree layers comprising the hyper-tree. 
\end{itemize}
To ensure symmetrical construction across all internal boundaries, our search pipeline strictly mandates that the total height $h$ must be perfectly divisible by the layer count $d$. 
Consequently, each individual XMSS tree layer operates with a uniform height denoted as $h'$:
\begin{equation*}
    h' = \frac{h}{d}
\end{equation*}
Adjusting the ratio between $h$ and $d$ governs the balance between signature size and key generation speed. 
Minimizing layers ($d$) reduces the total number of intermediate public keys and authentication paths included in the signature payload, saving valuable bytes. 
However, this exponentially increases the height $h'$ of each constituent tree layer, requiring significantly more hash operations to compute the root node during key generation.

\paragraph{Few-Time Signature Layer ($k$ and $a$)}
The lowest layer of the hyper-tree authenticates the Forest of Random Subsets (FORS) few-time signature scheme, which directly signs the message digest.
\begin{itemize}
    \item $k$ represents the number of independent, parallel Merkle trees operating within the FORS layer.
    \item $a$ defines the height of each individual FORS tree. This parameter dictates that each constituent tree contains exactly $t = 2^a$ leaves.
\end{itemize}

Increasing $k$ and $a$ strengthens the scheme against multi-target subset forgery attacks, allowing the structure to safely support larger capacities. 
However, expanding these dimensions directly increases the byte footprint of the appended authentication paths.

\paragraph{One-Time Signature Base ($w$)}
Internal tree roots are signed using the Winternitz One-Time Signature Plus (WOTS+) scheme. 
The base parameter $w$ defines the underlying digit representation used to split the message digest during signing. 
Our framework evaluates candidate configurations across standard baseline settings where $w \in \{16, 32, 256\}$.

Selecting a larger base value significantly shortens the required number of parallel hash chains ($l$), which reduces the overall WOTS+ signature payload size. 
This spatial efficiency trades against local compute cycles: larger base parameters require significantly longer individual hash chains, resulting in higher execution costs during public key generation and signature creation.

We set the ranges as follows: $(h,d,k,a,w) \in \{40,...,52\}\times\{2,...,26\}\times\{2,...,25\}\times\{8,...,21\}\times\{16,32,256\}$.